\notechapter{N-20250301}{Projective Schemes and Gluing: Building Projective Space from Affine Windows}

\section{Projective geometry returns with functions attached}

Earlier projective geometry used homogeneous coordinates to add points at infinity and remove affine exceptions. Scheme theory adds another layer: every affine window comes with a ring of functions, and the windows must be glued compatibly.

The construction that does this is \(\operatorname{Proj}\). It is to graded rings what \(\operatorname{Spec}\) is to ordinary rings.

\section{The projective line as the model}

The best first example is \(\mathbb P^1\). Take two affine lines:
\[
  U_0=\operatorname{Spec}R[t],
  \qquad
  U_1=\operatorname{Spec}R[s].
\]
On the region where both coordinates are nonzero, glue them by
\[
  s=\frac1t.
\]
The result is the projective line. The point is not the formula itself, but the construction style:
\[
  \text{projective space is made by gluing affine charts.}
\]

\section{Graded rings and homogeneous data}

Let
\[
  S=S_0\oplus S_1\oplus S_2\oplus\cdots
\]
be a graded ring. Projective geometry only respects homogeneous equations, because
\[
  [x_0:\cdots:x_n]=[\lambda x_0:\cdots:\lambda x_n].
\]
The construction \(\operatorname{Proj}S\) uses homogeneous prime ideals not containing the irrelevant ideal.

This definition sounds heavier than the picture, but it encodes a familiar rule: equations in projective space must be homogeneous.

\section{Projective space as Proj}

The scheme-theoretic projective space is
\[
  \mathbb P^n_R=\operatorname{Proj}R[x_0,\ldots,x_n],
\]
where the variables have degree \(1\).

The standard affine open where \(x_i\ne0\) is written
\[
  D_+(x_i).
\]
On it, the ratios
\[
  \frac{x_j}{x_i}
\]
serve as affine coordinates.

Thus \(\mathbb P^n\) is covered by \(n+1\) affine schemes. Projective geometry is global, but it is computed through affine windows.

\section{The projective plane through three windows}

For \(\mathbb P^2\), the standard opens are
\[
  D_+(x),
  \qquad
  D_+(y),
  \qquad
  D_+(z).
\]
On \(D_+(z)\), use
\[
  X=\frac{x}{z},
  \qquad
  Y=\frac{y}{z}.
\]
On \(D_+(x)\), use
\[
  U=\frac{y}{x},
  \qquad
  V=\frac{z}{x}.
\]
Each chart is affine. The projective plane is the result of gluing these affine pieces by ratio changes on overlaps.

This is the practical meaning of projective schemes: use affine algebra locally, then glue.

\section{A conic across charts}

Consider the projective conic
\[
  xz-y^2=0.
\]
On the chart \(z\ne0\), set \(X=x/z\) and \(Y=y/z\). The equation becomes
\[
  X=Y^2.
\]
On the chart \(x\ne0\), set \(U=y/x\) and \(V=z/x\). The equation becomes
\[
  V=U^2.
\]
The same projective curve has multiple affine faces. A single affine chart may hide what happens at infinity, but the projective object keeps all charts at once.

\section{Projective closure with scheme memory}

Given an affine equation \(f(x,y)=0\), homogenizing gives a projective equation
\[
  F(x,y,z)=0.
\]
Classically, this adds points at infinity. Scheme-theoretically, it also keeps track of nilpotent or embedded structure if such structure is present.

So \(\operatorname{Proj}\) is not merely a notation for old projective varieties. It is the projective version of the local-function viewpoint developed in scheme theory.

\section{What Proj adds}

The earlier projective note emphasized incidence and transformations. The \(\operatorname{Proj}\) viewpoint emphasizes functions and gluing.

\[
\begin{array}{c|c}
\text{homogeneous coordinates} & \text{homogeneous primes}\\
\text{affine charts} & \text{standard affine opens}\\
\text{projective equations} & \text{homogeneous ideals}\\
\text{points at infinity} & \text{glued affine schemes}
\end{array}
\]
These are not competing languages. They are layers of the same geometry.


\section{Why the irrelevant ideal is excluded}

For \(S=R[x_0,\ldots,x_n]\), the ideal
\[
  S_+=(x_0,\ldots,x_n)
\]
is called the irrelevant ideal.  It corresponds to the forbidden point where all homogeneous coordinates vanish at once:
\[
  [0:\cdots:0],
\]
which is not a projective point.

Thus \(\operatorname{Proj}S\) uses homogeneous prime ideals that do not contain all positive-degree elements.  This technical-looking exclusion is simply the algebraic version of the rule that homogeneous coordinates cannot all be zero.

\section{Standard opens as spectra}

For a homogeneous element \(f\in S\), the standard open \(D_+(f)\) is affine.  Its coordinate ring is the degree-zero part of the localization:
\[
  \mathcal O(D_+(f))=(S_f)_0.
\]
This formula explains why ratios appear in projective coordinates.  Localizing at \(x_i\) allows division by powers of \(x_i\), and taking degree zero keeps only scale-invariant expressions such as
\[
  \frac{x_j}{x_i}.
\]
So affine coordinates on projective charts are not arbitrary.  They are degree-zero fractions.

\section{A second conic chart calculation}

For the conic
\[
  xz-y^2=0,
\]
the chart \(y\ne0\) uses coordinates
\[
  A=\frac{x}{y},
  \qquad
  B=\frac{z}{y}.
\]
The equation becomes
\[
  AB=1.
\]
So in this chart the conic looks like a hyperbola.  In another chart it looked like a parabola.  This is a concrete reminder that affine appearances depend on the chosen window.

The projective curve is the object obtained by holding all these windows together.

\section{Line bundles are already nearby}

Projective geometry naturally leads to line bundles.  On \(\mathbb P^n\), the twisting sheaves \(\mathcal O(d)\) organize homogeneous polynomials of degree \(d\).  A homogeneous equation of degree \(d\) is better thought of as a section of \(\mathcal O(d)\) than as an ordinary global function.

This explains why projective varieties have fewer global regular functions than affine varieties.  Projective geometry shifts attention from functions to sections of line bundles.

This remark points back to divisors and Riemann--Roch: projective geometry, line bundles, and spaces of sections are parts of the same story.

The twisting sheaves therefore complete the passage from homogeneous coordinates to projective schemes.  The same graded algebra that builds \(\operatorname{Proj}S\) also supplies its affine windows and its line bundles: projective geometry has been reconstructed from homogeneous algebra and gluing while retaining the local functions invisible in a purely incidence-based picture.
